% GENERATED FILE. Edit canonical lessons, then run npm run generate:books.
% canonical-source-hash: fnv1a64:0fab715883ba221b
% canonical-lessons: ZH-C112-chukou, ZH-W112-chu, ZH-C112-rukou, ZH-C112-tui, ZH-C112-la, ZH-C112-louti

\chapter{Way In, Way Out}
\label{ch:zh-way-in-way-out}
\hlchaptermodality{112}

\begin{chapteropening}
\textbf{By the end of this chapter:} \emph{I can say the words on signs for the way out, the way in, push, pull and the stairs.}

Complete the last lesson of chapter 112: I can say the words on signs for the way out, the way in, push, pull and the stairs.
\end{chapteropening}

\section[\texorpdfstring{\zh{出口} (chūkǒu)}{出口 (chūkǒu)}]{\zh{出口} (chūkǒu) — exit; way out}

\label{lesson:ZH-C112-chukou}



\begin{quote}
Before the new one: say the Mandarin for crossroads, then the Mandarin for to turn a corner.
\end{quote}

\subsection*{You'll want to know: \zh{出口}}
\textbf{\zh{出口}} — \emph{chūkǒu} — \textquotedblleft{}exit; way out\textquotedblright{}.

\begin{quote}
\textbf{\zh{出口在前边}} \emph{chūkǒu zài qiánbian} — The way out is up ahead.
\end{quote}

The new character, in pieces: \zh{出} — two halves like \zh{山}, one on top of the other, with one long vertical down through both. The next lesson puts it in your hand.

\begin{cousinweb}
\textbf{\zh{出口}} (\emph{chūkǒu}) is built from \zh{出} 'to go out' and \zh{口} 'mouth, opening', both inherited from Middle Chinese, \emph{tsyhwit} and \emph{khuwX} (Guangyun, in Baxter's transcription): the opening you go out by.
\end{cousinweb}

\begin{grammarlens}[title={What you've built}]
One more word for signs and notices.
\end{grammarlens}

\subsection*{Your turn}
\begin{itemize}
\raggedright
  \item \emph{Say it:} \emph{chūkǒu}
  \item \emph{Say it:} \emph{chūkǒu}, once more
  \item \emph{Say it:} \emph{chūkǒu zài qiánbian} — The way out is up ahead.
  \item \emph{Recall:} say the Mandarin for crossroads, then the Mandarin for to turn a corner, then say \emph{chūkǒu} again
\end{itemize}

\begin{tcolorbox}[breakable,colback=teal!4,colframe=teal!35!black,title={Before you move on}]
What does \zh{出口} mean? (\textquotedblleft{}Exit; way out\textquotedblright{}.) Say it once more.
\end{tcolorbox}
\section[\texorpdfstring{\zh{出} (chū)}{出 (chū)}]{\zh{出} — two halves like \zh{山}, one long vertical}

\label{lesson:ZH-W112-chu}



\begin{quote}
The word first, then the character inside it.

\begin{itemize}
\raggedright
  \item \emph{Recall:} say \emph{chūkǒu} — \textbf{R1}, one lesson back, before it decays
\end{itemize}

\begin{quote}
\zh{山} \zh{口}
\end{quote}

\emph{Read it:} each one aloud, with nothing said first — \emph{shān}, \emph{kǒu}
\end{quote}

\subsection*{Script — two halves like \zh{山}, one long vertical}
\hlblockfigure{\detokenize{figures/ZH-W112-chu-filmstrip.pdf}}{How \zh{出} is written, stroke by stroke}

\begin{quote}
\zh{出}
\end{quote}

\textbf{The top half, two strokes}: a short vertical that turns right along its foot, then the short vertical on the right.

\textbf{The long vertical}, down the middle, through the top half and on into the bottom one.

\textbf{The bottom half, two strokes}, wider: a vertical that turns right along the foot, then the vertical on the right.

2 + 1 + 2 = five. Unlike \zh{山}, the middle vertical comes third.

\subsection*{Your turn}
\begin{itemize}
\raggedright
  \item \emph{Write it:} \zh{出}, all five, with the model in front of you
  \item \emph{Cover:} write it again from memory, and count the pieces, not the strokes
\end{itemize}

\begin{tcolorbox}[breakable,colback=teal!4,colframe=teal!35!black,title={Before you move on}]
How many strokes in \zh{出}? (\textbf{Five}.) Which stroke comes third? (\textbf{The long vertical down the middle}.) Which character you know does it look like, twice? (\textbf{\zh{山}}.)

Source: \href{https://raw.githubusercontent.com/chanind/hanzi-writer-data/68d10a4b21150cae5e1ebbd223eed289cf32d90c/data/\%E5\%87\%BA.json}{Hanzi Writer Data \zh{出}, pinned stroke paths and medians}.
\end{tcolorbox}
\section[rùkǒu]{rùkǒu — entrance; way in}

\label{lesson:ZH-C112-rukou}



\begin{quote}
Before the new one: say the Mandarin for to turn a corner, then the Mandarin for exit.
\end{quote}

\subsection*{You'll want to know: rùkǒu}
\emph{rùkǒu} — \textquotedblleft{}entrance; way in\textquotedblright{}.

\begin{quote}
\emph{rùkǒu zài nàr} — The way in is over there.
\end{quote}

Its characters come later; for now it is a word for the ear.

\begin{cousinweb}
\emph{rùkǒu} is built from rù 'to go in' and \zh{口} 'opening', both inherited from Middle Chinese, \emph{nyip} and \emph{khuwX} (Guangyun, in Baxter's transcription), glossed 'inside, to take in' and 'mouth'.
\end{cousinweb}

\begin{grammarlens}[title={What you've built}]
One more word for signs and notices.
\end{grammarlens}

\subsection*{Your turn}
\begin{itemize}
\raggedright
  \item \emph{Say it:} \emph{rùkǒu}
  \item \emph{Say it:} \emph{rùkǒu}, once more
  \item \emph{Say it:} \emph{rùkǒu zài nàr} — The way in is over there.
  \item \emph{Recall:} say the Mandarin for to turn a corner, then the Mandarin for exit, then say \emph{rùkǒu} again
\end{itemize}

\begin{tcolorbox}[breakable,colback=teal!4,colframe=teal!35!black,title={Before you move on}]
What does rùkǒu mean? (\textquotedblleft{}Entrance; way in\textquotedblright{}.) Say it once more.
\end{tcolorbox}
\section[tuī]{tuī — to push}

\label{lesson:ZH-C112-tui}



\begin{quote}
Before the new one: say the Mandarin for exit, then the Mandarin for entrance.
\end{quote}

\subsection*{You'll want to know: tuī}
\emph{tuī} — \textquotedblleft{}to push\textquotedblright{}.

\begin{quote}
\emph{qǐng tuī mén} — Push the door, please.
\end{quote}

Its characters come later; for now it is a word for the ear.

\begin{cousinweb}
\emph{tuī} is inherited from Middle Chinese \emph{thwoj} (Guangyun, in Baxter's transcription). The dictionary glosses the character 'to push' under a second reading, \emph{tsyhwij}, and lists \emph{thwoj} beside it; tuī continues \emph{thwoj}.
\end{cousinweb}

\begin{grammarlens}[title={What you've built}]
One more word for signs and notices.
\end{grammarlens}

\subsection*{Your turn}
\begin{itemize}
\raggedright
  \item \emph{Say it:} \emph{tuī}
  \item \emph{Say it:} \emph{tuī}, once more
  \item \emph{Say it:} \emph{qǐng tuī mén} — Push the door, please.
  \item \emph{Recall:} say the Mandarin for exit, then the Mandarin for entrance, then say \emph{tuī} again
\end{itemize}

\begin{tcolorbox}[breakable,colback=teal!4,colframe=teal!35!black,title={Before you move on}]
What does tuī mean? (\textquotedblleft{}To push\textquotedblright{}.) Say it once more.
\end{tcolorbox}
\section[lā]{lā — to pull}

\label{lesson:ZH-C112-la}



\begin{quote}
Before the new one: say the Mandarin for entrance, then the Mandarin for to push.
\end{quote}

\subsection*{You'll want to know: lā}
\emph{lā} — \textquotedblleft{}to pull\textquotedblright{}.

\begin{quote}
\emph{bú yào tuī qǐng lā} — Don't push, pull.
\end{quote}

Its characters come later; for now it is a word for the ear.

\begin{cousinweb}
\emph{lā} is inherited from Middle Chinese \emph{lop} (Guangyun, in Baxter's transcription), glossed 'to break' and 'to crush'. The regular change from \emph{lop} would give là, not lā, and pulling is not given there; how the word came to its sound and sense is not settled.
\end{cousinweb}

\begin{grammarlens}[title={What you've built}]
One more word for signs and notices.
\end{grammarlens}

\subsection*{Your turn}
\begin{itemize}
\raggedright
  \item \emph{Say it:} \emph{lā}
  \item \emph{Say it:} \emph{lā}, once more
  \item \emph{Say it:} \emph{bú yào tuī qǐng lā} — Don't push, pull.
  \item \emph{Recall:} say the Mandarin for entrance, then the Mandarin for to push, then say \emph{lā} again
\end{itemize}

\begin{tcolorbox}[breakable,colback=teal!4,colframe=teal!35!black,title={Before you move on}]
What does lā mean? (\textquotedblleft{}To pull\textquotedblright{}.) Say it once more.
\end{tcolorbox}
\section[lóutī]{lóutī — stairs; staircase}

\label{lesson:ZH-C112-louti}



\begin{quote}
Before the new one: say the Mandarin for to push, then the Mandarin for to pull.
\end{quote}

\subsection*{You'll want to know: lóutī}
\emph{lóutī} — \textquotedblleft{}stairs; staircase\textquotedblright{}.

\begin{quote}
\emph{lóutī zài ménkǒu pángbiān} — The stairs are next to the doorway.
\end{quote}

Its characters come later; for now it is a word for the ear.

\begin{cousinweb}
\emph{lóutī} is built from lóu and tī, both inherited from Middle Chinese, \emph{luw} and \emph{thej} (Guangyun, in Baxter's transcription), glossed 'a house built on a house' and 'wooden steps'.
\end{cousinweb}

\begin{grammarlens}[title={What you've built}]
One more word for signs and notices.
\end{grammarlens}

\subsection*{Your turn}
\begin{itemize}
\raggedright
  \item \emph{Say it:} \emph{lóutī}
  \item \emph{Say it:} \emph{lóutī}, once more
  \item \emph{Say it:} \emph{lóutī zài ménkǒu pángbiān} — The stairs are next to the doorway.
  \item \emph{Recall:} say the Mandarin for to push, then the Mandarin for to pull, then say \emph{lóutī} again
\end{itemize}

\begin{tcolorbox}[breakable,colback=teal!4,colframe=teal!35!black,title={Before you move on}]
What does lóutī mean? (\textquotedblleft{}Stairs; staircase\textquotedblright{}.) Say it once more.
\end{tcolorbox}
